% !TEX program = tectonic
\documentclass[11pt,a4paper]{article}
\usepackage[margin=1.5cm]{geometry}
\usepackage[T1]{fontenc}
\usepackage[utf8]{inputenc}
\usepackage{lmodern}
\usepackage[french]{babel}
\usepackage{amsmath,amssymb}
\usepackage{xcolor}
\usepackage{booktabs}
\usepackage{enumitem}
\usepackage{mdframed}
\definecolor{bleu}{RGB}{20,77,144}
\definecolor{vert}{RGB}{20,110,60}
\definecolor{rouge}{RGB}{160,30,30}
\definecolor{grisf}{RGB}{100,100,100}
\pagestyle{empty}
\setlength{\parindent}{0pt}
\setlength{\parskip}{2pt}

\newcommand{\rem}[2]{\vspace{2mm}{\color{bleu}\large\bfseries #1 --- #2}\vspace{0.5mm}\par}
\newcommand{\cit}[1]{\begin{mdframed}[linecolor=grisf!50, backgroundcolor=grisf!7, linewidth=0.6pt,
  roundcorner=3pt, innertopmargin=3pt, innerbottommargin=3pt]\small\itshape \og #1\fg{}\end{mdframed}}
\newcommand{\trad}[1]{{\color{bleu}\small\textbf{Traduction :} #1}\par\vspace{0.6mm}}
\newcommand{\prop}[2]{\vspace{0.6mm}\begin{mdframed}[linecolor=vert!60, backgroundcolor=vert!5, linewidth=0.6pt,
  roundcorner=3pt, innertopmargin=3pt, innerbottommargin=3pt]{\color{vert}\textbf{Propriété #1.}} #2\end{mdframed}\vspace{0.4mm}}
\newcommand{\mes}[1]{{\color{rouge}\small$\blacktriangleright$ \textbf{mesuré :} #1}\par}
\newcommand{\niche}[1]{{\color{vert}\small$\hookrightarrow$ \textbf{redonne l'existant :} #1}\par}

\begin{document}

\begin{center}
{\LARGE\bfseries Les neuf remarques --- et la mathématique qui y répond}\\[1mm]
{\normalsize backtest \textbf{par membre} $\cdot$ une section par remarque : \emph{la remarque traduite en énoncé, les objets, les propriétés, la mesure}}\\[0.5mm]
{\small\color{grisf} remarques d'Olivier Herbout sur \texttt{SLIDES\_DONNEES\_S1S2\_V2} --- 3 dans son message (R), 6 annotées dans le PDF (A)}
\end{center}

\vspace{0.5mm}
\begin{mdframed}[linecolor=grisf, backgroundcolor=grisf!6, linewidth=0.8pt, roundcorner=4pt]
\textbf{Notations communes.} $\mathcal{C}$ = calendrier de bourse ; $\tau_k$ = date de transaction, $\delta_k$ = date de
\textbf{divulgation}, $d_k \in \{-1,+1\}$ = sens, $A_k$ = montant (milieu de fourchette), $m_k$ = élu, $i_k$ = titre.
$n_i(t)$ = parts détenues, $P_i(t)$ = prix, $w_i(t)$ = poids, $r_P(t)$ = rendement du portefeuille,
$r_M$ = marché, $r_f$ = taux sans risque.
\textbf{Toute nouveauté est un opérateur à un paramètre dont un réglage redonne exactement l'existant} :

\vspace{1.5mm}
\begin{center}\footnotesize
\begin{tabular}{@{}l l c l@{}}
\toprule
\textbf{remarque} & \textbf{objet mathématique} & \textbf{défaut} & \textbf{redonne l'existant} \\
\midrule
R1 $\cdot$ A3 & noyau de décote $\varphi_H(a)$ sur l'âge des lots & $H = 252$ j & $\varphi \equiv 1$ \\
A1 & poids de rafale $\psi(B) = 1/B$, décompte $n_{\text{eff}}$ & --- & $\psi \equiv 1$ \\
R2 & excès signé $\mathrm{AR}_k$ $+$ budget de vente $\theta$ & $\theta$ balayé & $\theta = 0$ \\
R3 $\cdot$ A5 $\cdot$ A6 & $\beta$ annuel $\cdot$ Carhart-4 $\cdot$ excès $r_f$ $\cdot$ NW & $L = 21$ & 1 bloc, 1 facteur, sans $r_f$ \\
A2 & filtre d'aller-retour $\Lambda_{x,\varepsilon}$ & $\varepsilon = 10\,\%$ & $\varepsilon \to \infty$ : clip brut \\
A4 & rétrécissement $\to$ classement a posteriori & $\hat\tau^2$ (moments) & $\tau^2 \to \infty$ $+$ la porte \\
\bottomrule
\end{tabular}
\end{center}
\end{mdframed}

% ═══════════════════════════════════════════════════════════
\rem{R1 $\cdot$ A3}{Les positions dormantes --- le noyau de décote}
\cit{si les holdings sont stables pendant trop longtemps (à définir), ce ne peut être considéré comme de
l'inside trading : il faut les discounter // les supprimer (le plus propre, les discounter et les faire tendre vers 0)}
\vspace{-1.5mm}
\cit{faire une analyse de sensibilité : retirer du portefeuille les trades qui sont là depuis plus de 1 an, plus de 2 ans, 3 ans, 5 ans}

\trad{le portefeuille observé mélange deux composantes --- des décisions récentes (informatives) et de
l'inertie (non informative). On cherche l'\textbf{opérateur qui projette sur la composante récente}, indexé par
une échelle de temps $H$, et tel que $H \to \infty$ redonne le portefeuille entier.}

\textbf{Les objets.} L'appariement FIFO définit à chaque date $t$ l'ensemble des \textbf{lots ouverts} du titre $i$ :
$\mathcal{L}_i(t) = \{(\delta_\ell, q_\ell)\}$, où $q_\ell$ est la part restante du lot $\ell$ et $\delta_\ell$ sa date de
divulgation. L'âge se compte en jours de bourse depuis $\delta_\ell$ --- \emph{on date l'information, pas l'opération} :
\[
a_\ell(t) = \#\bigl(\mathcal{C} \cap (\delta_\ell,\, t]\bigr),
\qquad
\varphi_H(a) = 2^{-a/H},
\qquad
n_i^{\varphi}(t) = \sum_{\ell \in \mathcal{L}_i(t)} \varphi_H\bigl(a_\ell(t)\bigr)\, q_\ell
\]
\[
w_i^{\varphi}(t) = \frac{n_i^{\varphi}(t)\, P_i(t)}{\sum_j n_j^{\varphi}(t)\, P_j(t)},
\qquad
r_P^{\varphi}(t) = \sum_i w_i^{\varphi}(t{-}1)\, r_i(t),
\qquad
\bar A^{\varphi}(t) = \frac{\sum_i \sum_{\ell \in \mathcal{L}_i(t)} \varphi_H(a_\ell)\, q_\ell P_i\, a_\ell(t)}{\sum_i \sum_{\ell} \varphi_H(a_\ell)\, q_\ell P_i}
\]

\prop{1 (nichage)}{$\varphi \equiv 1$ (soit $H \to \infty$) donne $n_i^{\varphi} = n_i$, donc $w^{\varphi} = w$ et
$r_P^{\varphi} = r_P$ : \textbf{le moteur actuel est le cas limite}. Et $\varphi_H(a) = \mathbf{1}\{a \le H\}$ avec
$H \in \{1,2,3,5\}$ ans produit \emph{exactement} l'analyse de sensibilité demandée --- la troncature dure et la
décote douce sont deux réglages du même opérateur.}

\prop{2 (le portefeuille reste un portefeuille)}{$\sum_i w_i^{\varphi}(t) = 1$ par construction, tant qu'un lot
survit. $r_P^{\varphi}$ est donc une moyenne pondérée de rendements de titres : c'est un vrai rendement de
portefeuille, autofinancé. \emph{C'est la renormalisation qui l'assure} --- décoter les parts sans renormaliser
ferait fuir la valeur et mélangerait décote et performance.}

\prop{3 (mémoire effective)}{$\sum_{a \ge 0} \varphi_H(a) = \bigl(1 - 2^{-1/H}\bigr)^{-1} \approx H/\ln 2 \approx 1{,}44\,H$.
Une demi-vie de 252 j correspond donc à une mémoire effective de $\approx 363$ jours de bourse : le portefeuille
\og se souvient \fg{} d'environ un an et demi d'activité.}

\textbf{Le constat.} FIFO en dollars par (membre, titre), lots ouverts à chaque 31/12 :

\begin{center}\footnotesize
\begin{tabular}{@{}l r r r r@{}}
\toprule
coupe & $>1$ an & $>2$ ans & $>3$ ans & âge moyen pondéré \\
\midrule
2016 & 67,9 \% & 16,0 \% & \phantom{0}1,8 \% & 1,5 an \\
2020 & 79,2 \% & 63,0 \% & 47,4 \% & 3,3 ans \\
2025 & 91,6 \% & 87,5 \% & 83,8 \% & 6,1 ans \\
\midrule
\textbf{moyenne des 10 coupes} & \textbf{79,1 \%} & \textbf{60,6 \%} & \textbf{45,3 \%} & \textbf{3,4 ans} \\
\bottomrule
\end{tabular}
\end{center}
\vspace{-1mm}
{\footnotesize\color{grisf} les coupes 2016--2018 sont censurées à gauche (la donnée commence en 2014) ; le régime établi est 2023--2025, où le $>1$ an vaut 91 à 94 \%.}

\mes{part de la valeur qui survit à la décote : $H = 6$ mois $\to 13\,\%$ $\cdot$ $H = 1$ an $\to \mathbf{25{,}4\,\%}$ $\cdot$ $H = 3$ ans $\to 46\,\%$. \textbf{À près de 80 \% en valeur, ce qu'on mesure est du dormant.}}

{\small\textbf{Le choix de $H$, fixé d'avance :} la durée médiane d'un épisode achat$\to$revente vaut \textbf{349 j} sur nos
données. Poser $H = 252$ j place la demi-vie au voisinage du comportement observé, sans l'ajuster sur les
rendements. \textbf{Un seul paramètre, jamais recalibré.}}

% ═══════════════════════════════════════════════════════════
\rem{A1}{La délégation --- une ligne déclarée n'est pas une décision}
\cit{quand on délègue c'est notamment pour ne pas faire de l'inside trading. A supprimer de l'analyse assez rapidement}

\trad{le STOCK Act compte des \textbf{lignes}, pas des \textbf{décisions}. Un gérant qui rebalance un portefeuille émet
$B$ lignes le même jour, qui portent \emph{un seul} geste. La bonne réponse n'est pas de deviner qui délègue, mais
de \textbf{corriger l'unité d'observation} : pondérer chaque ligne par $1/B$.}

\textbf{Les objets.} $B_{m,d}$ = nombre de lignes déclarées par le membre $m$ le jour $d$ ; poids et décompte effectif :
\[
B_{m,d} \;=\; \#\bigl\{\, k \;:\; m_k = m,\ \tau_k = d \,\bigr\},
\qquad
\psi(B) = \frac{1}{B},
\qquad
n_{\text{eff}}(m) \;=\; \sum_{k \,:\, m_k = m} \psi\bigl(B_{m,\tau_k}\bigr)
\]

\prop{4 (le poids $1/B$ compte exactement les jours de déclaration)}{Pour un membre $m$ et un jour $d$ où il
déclare $B_{m,d}$ lignes, la contribution de ce jour vaut $\sum_{k} 1/B_{m,d} = 1$. En sommant,
$n_{\text{eff}}(m) = \#\{\text{jours où } m \text{ déclare}\}$, \textbf{exactement}. Le poids $1/B$ est donc le seul qui
rende la somme égale au nombre de \emph{gestes datés}, et non au nombre de lignes. Aucun paramètre libre ;
$\psi \equiv 1$ redonne le décompte en lignes.
\emph{(Vérifié numériquement : $\sum_k \psi_k = 24\,820{,}000000$ pour 24\,820 jours-membre, écart nul.)}}

\begin{center}\footnotesize
\begin{tabular}{@{}l r r r@{}}
\toprule
membre & lignes $n$ & \textbf{jours-membre} $n_{\text{eff}}$ & $n_{\text{eff}}/n$ \\
\midrule
Khanna & 35\,046 & 1\,779 & \textbf{5,1 \%} \\
McCaul & 20\,296 & 2\,052 & 10,1 \% \\
Whitehouse & \phantom{0,0}763 & \phantom{0,0}219 & 28,7 \% \\
\textbf{Pelosi} & \phantom{0,0}141 & \phantom{0,00}77 & \textbf{54,6 \%} \\
\bottomrule
\end{tabular}
\end{center}
\vspace{-1mm}
\mes{la population passe de \textbf{134\,464 lignes} à \textbf{24\,820 jours-membre} (18,5 \%). Les deux plus gros
déposants pèsent \textbf{41,2 \% des lignes} mais seulement \textbf{15,4 \% des jours} : leur poids est divisé par 2,7
\emph{sans exclure personne à la main}. L'éligibilité en jours plutôt qu'en lignes fait passer les membres jugeables
de \textbf{264} à \textbf{199}.}

\begin{mdframed}[linecolor=rouge!60, backgroundcolor=rouge!5, linewidth=0.6pt, roundcorner=3pt,
  innertopmargin=3pt, innerbottommargin=3pt]
{\small\textbf{Ce que cet objet ne fait PAS --- et qu'il ne faut pas lui faire dire.} Il \emph{ne détecte pas} la
délégation. Testé : la rafale médiane de McCaul le place au 91\ieme{} centile, pas en dehors de la distribution ;
et la statistique dépend de l'agrégation choisie (chez Khanna, rafale médiane \textbf{133} vue depuis un trade,
mais \textbf{2} vue depuis un jour --- un facteur 60). \textbf{Aucun seuil de rafale ne sépare proprement les gérants
des élus.} Ce que $1/B$ règle, c'est autre chose et c'est solide : \emph{un rebalancement de trust cesse de compter
pour 133 observations indépendantes}. On corrige le décompte, on ne prétend pas lire l'intention.}
\end{mdframed}

% ═══════════════════════════════════════════════════════════
\rem{R2}{Achat contre vente --- l'information est-elle symétrique en son signe ?}
\cit{j'aurais envie de savoir s'il y a une différence d'information // de valeur entre les long et les short positions}

\trad{nos données ne contiennent pas de vraies positions courtes (un PTR déclare des achats et des ventes de
titres détenus). La question devient : $\mathbb{E}[\text{information} \mid d_k = +1] = \mathbb{E}[\text{information} \mid d_k = -1]$ ?
Trois niveaux de réponse : l'événement, le portefeuille, et l'asymétrie \emph{a priori} du choix.}

\textbf{(a) Au niveau de l'événement.} Excès signé, ancré sur la divulgation, puis agrégé en cohortes :
\[
\mathrm{AR}_k(H) = d_k\Bigl[ R_{i_k}\bigl(\delta_k \!\to\! \delta_k{+}H\bigr) - R_{\mathrm{ref}}\bigl(\delta_k \!\to\! \delta_k{+}H\bigr) \Bigr],
\qquad
\overline{\mathrm{AR}}_\mu = \frac{1}{|\mathcal{K}_\mu|}\sum_{k \in \mathcal{K}_\mu} \mathrm{AR}_k(H)
\]
où $\mathcal{K}_\mu$ = les trades divulgués le mois $\mu$. Le test porte sur les $T$ observations mensuelles :
$t = \overline{\mathrm{AR}}/(\hat\sigma/\sqrt{T})$, et $H_0 : \mathbb{E}[\mathrm{AR}\mid{+}1] = \mathbb{E}[\mathrm{AR}\mid{-}1]$.
\textbf{L'agrégation n'est pas cosmétique} : deux trades séparés de moins de $H$ jours partagent leur fenêtre, donc
$\mathrm{Cov}(\mathrm{AR}_k, \mathrm{AR}_j) \ne 0$ ; le $t$ calculé sur les trades individuels sous-estime l'écart-type
d'un facteur croissant avec $H$.

\textbf{(b) Au niveau du portefeuille}, avec un budget de vente $\theta \in [0,1]$ : $r_\theta = r_L - \theta\, r_S$.

\prop{5 (le long-short achète de la variance, pas de l'alpha)}{Les MCO étant linéaires en la variable expliquée,
sur les mêmes facteurs :
\[
\alpha_\theta = \alpha_L - \theta\,\alpha_S
\qquad\text{mais}\qquad
\sigma^2_\theta = \sigma_L^2 + \theta^2 \sigma_S^2 - 2\theta\,\sigma_{LS},
\qquad
t_\theta = \frac{\alpha_L - \theta \alpha_S}{\sigma_\theta}\sqrt{n}
\]
$\alpha$ est \emph{linéaire} en $\theta$, $t$ ne l'est pas. Le minimum de variance est atteint en
$\theta^\star = \sigma_{LS}/\sigma_S^2$. \textbf{Donc la jambe vendeuse ne peut aider que par la variance résiduelle} ---
et $\theta = 0$ redonne exactement le portefeuille long actuel.}

\textbf{(c) L'asymétrie a priori du choix.} Sous $H_0$ \og aucune information \fg{}, un achat désigne un titre dans
l'univers investissable $\mathcal{U}_t$, une vente dans le seul portefeuille détenu $\mathcal{H}_m(t)$ :
\[
\mathcal{I}_{\text{achat}} = \log_2 |\mathcal{U}_t| \approx \log_2 1027 = 10{,}0 \text{ bits},
\qquad
\mathcal{I}_{\text{vente}} = \log_2 |\mathcal{H}_m(t)| \approx \log_2 14 = 3{,}8 \text{ bits}
\]
soit $\Delta \approx \mathbf{6{,}2}$ \textbf{bits} : un achat désigne un titre parmi mille, une vente parmi quatorze.
\emph{C'est une asymétrie de sélectivité, pas une prédiction de rendement --- à ne pas sur-interpréter.}

\mes{la part de ventes atteint \textbf{56,1 \%} en décembre contre 48,7 \% de janvier à novembre ($z = +13{,}7$) : la
récolte de moins-values fiscales est un motif \textbf{non informationnel, massif et daté}, qui pollue la jambe
vendeuse et qu'aucun modèle de rendement ne peut absorber.}

{\small\textbf{Le biais à déclarer} : le moteur tronque les ventes par $q^- = \min(q, h)$ --- \textbf{937 M\$} de masse
vendue non couverte dans le run témoin --- et 44 \% des positions ne sont jamais revendues. \textbf{L'échantillon des
ventes est censuré}, toute comparaison des deux sens en hérite ; à défaut de correction, on publie la borne.}

% ═══════════════════════════════════════════════════════════
\rem{R3 $\cdot$ A5 $\cdot$ A6}{Le $\beta$ --- quatre corrections, et d'abord une réponse}
\cit{le calcul du beta est un peu trop naïf} \vspace{-1.5mm}
\cit{bon test mais s'il est seul, est trop naïf. Le beta est il stable dans le temps par exemple ?}

\trad{le modèle actuel est $r_P = a + b\,r_M + e$ : \textbf{un} $\beta$, \textbf{un} facteur, \textbf{sans} $r_f$,
\textbf{sans} correction d'autocorrélation. Quatre défauts identifiables, quatre corrections emboîtées.}

\textbf{La réponse à sa question, d'abord.} $\beta$ ré-estimé année par année :

\begin{center}\scriptsize
\begin{tabular}{@{}l r l r@{}}
\toprule
membre & $\beta$ global & $\beta$ par année, 2016 $\to$ 2025 & amplitude \\
\midrule
\textbf{Whitehouse} {\color{grisf}(le plus solide du deck)} & 1,19 & 1,01 $\cdot$ 1,09 $\cdot$ 1,13 $\cdot$ 1,10 $\cdot$ 1,03 $\cdot$ 1,22 $\cdot$ 1,25 $\cdot$ 1,27 $\cdot$ \textbf{1,88} $\cdot$ 1,45 & \textbf{0,88} \\
Pelosi & 1,14 & 0,82 $\cdot$ 1,09 $\cdot$ 1,30 $\cdot$ 1,33 $\cdot$ 0,91 $\cdot$ 1,18 $\cdot$ 1,40 $\cdot$ 1,30 $\cdot$ \textbf{1,53} $\cdot$ 1,35 & \textbf{0,71} \\
Rooney & 0,55 & \phantom{0,00 $\cdot$} 0,79 $\cdot$ 0,48 $\cdot$ 0,51 $\cdot$ 0,64 $\cdot$ 0,58 $\cdot$ 0,53 $\cdot$ 0,55 $\cdot$ 0,44 $\cdot$ 0,43 & 0,36 \\
\bottomrule
\end{tabular}
\end{center}
\vspace{-1.5mm}
{\footnotesize\color{grisf} l'amplitude vaut 30 à 75 \% du niveau ; le $\beta$ de Whitehouse quasi double. Un $\beta$ unique sur douze ans n'a pas de sens --- et la stratégie top-$K$, elle, \emph{refait son portefeuille chaque année}.}

\prop{6 (le $\beta$ unique n'est pas une moyenne des années --- c'est celui des années agitées)}{Découpons
l'échantillon en blocs $g$ et posons $S_{xx,g} = \sum_{t \in g}(x_t - \bar x)^2$ avec $x = r_M - r_f$. Alors
\[
\hat\beta_{\text{global}} \;=\; \sum_g \varpi_g\, \hat\beta_g,
\qquad
\varpi_g = \frac{S_{xx,g}}{\sum_h S_{xx,h}}
\]
Les poids sont proportionnels à la \textbf{variance du marché} dans le bloc, pas à sa durée.
\mes{sur 2016--2025 : \textbf{2020 pèse 34,4 \%} et 2022 18,5 \% --- à eux deux \textbf{53 \%} --- tandis que 2017 pèse
\textbf{1,5 \%}. Le nombre effectif d'années vaut $1/\sum_g \varpi_g^2 = \mathbf{5{,}4}$ sur 10.
\textbf{Le $\beta$ de 1,20 affiché est, pour moitié, le $\beta$ de 2020 et 2022.}}}

\textbf{(i) Le marché en excès.} En régressant du brut sur du brut, on estime en fait
$a = \alpha + r_f(1-\beta)$, d'où :
\[
r_P - r_f = \alpha + \beta\,(r_M - r_f) + \varepsilon
\qquad\Longrightarrow\qquad
\alpha \;=\; a \;-\; r_f\,(1-\beta)
\]
\mes{$r_f$ vaut \textbf{1,87 \%/an} en moyenne (\textbf{4,85 \%} sur 2023--2026). Avec $\beta = 1{,}198$, le terme
$r_f(1-\beta) = -0{,}37\,\%$ : l'intercept brut \textbf{sous-estime} $\alpha$ de 0,37 \%/an (0,96 \% au régime récent).
\emph{Cette correction joue en notre faveur} --- il faut le dire avant de la faire.}

\textbf{(ii) Sortir du facteur unique} --- les facteurs sont déjà dans le cache :
\[
r_P - r_f = \alpha + \beta_M\,\mathrm{MKT} + \beta_S\,\mathrm{SMB} + \beta_H\,\mathrm{HML} + \beta_W\,\mathrm{WML} + \varepsilon
\]

\textbf{(iii) Tester la stabilité} --- $\beta$ interagi avec les $J$ blocs annuels, puis test de Chow :
\[
r_P - r_f = \alpha + \sum_{Y} \beta_Y\, \mathbf{1}\{t \in Y\}\,(r_M - r_f) + \varepsilon,
\qquad
F = \frac{(\mathrm{SCR}_c - \mathrm{SCR}_u)/(J-1)}{\mathrm{SCR}_u/(n-2J)} \;\sim\; F_{J-1,\; n-2J}
\]
sous $H_0 : \beta_{2015} = \cdots = \beta_{2025}$, où $\mathrm{SCR}_c$ et $\mathrm{SCR}_u$ sont les sommes des carrés
des résidus contrainte et non contrainte.

\textbf{(iv) Des écarts-types honnêtes} --- Newey-West, $L = 21$ (la convention du notebook 08) :
\[
\hat\sigma^2_{\mathrm{NW}} \;=\; \hat\gamma_0 + 2\sum_{j=1}^{L}\Bigl(1 - \tfrac{j}{L+1}\Bigr)\hat\gamma_j
\]
{\small\emph{Attention au sens} : on dit souvent que la correction HAC \og élargit \fg{} les écarts-types. Ce n'est vrai
que si l'autocorrélation est positive. Ici $\hat\rho_1$ de l'excès quotidien est \textbf{négatif}, donc
$\hat\sigma_{\mathrm{NW}} < \hat\sigma_{\mathrm{OLS}}$ et le $t$ \emph{monte} légèrement. \textbf{On ne peut pas annoncer
le signe à l'avance --- on le mesure.}}

\prop{7 (pourquoi un $\beta$ court est du bruit)}{$\;\mathrm{sd}(\hat\beta) = \sigma_\varepsilon / (\sigma_M \sqrt{n})$.
Pour un membre à $n = 85$ jours avec $\sigma_\varepsilon \approx \sigma_M$, on obtient
$\mathrm{sd}(\hat\beta) \approx 1/\sqrt{85} \approx 0{,}11$, soit un intervalle à deux écarts-types de
$\pm 0{,}22$. \textbf{Un $\beta$ de 1,20 estimé sur quatre mois est indiscernable de 1,00} --- et c'est ce $\beta$ qui
sert à retirer le marché.}

\niche{(i) sans $r_f$, (ii) sans SMB/HML/WML, (iii) un seul bloc, (iv) écarts-types i.i.d. --- c'est exactement le modèle actuel.}

{\small\textbf{À annoncer avant de mesurer :} les corrections poussent en sens \emph{opposés} --- $r_f$ remonte $\alpha$,
les facteurs de style le baissent (portefeuilles chargés en méga-capitalisations technologiques, donc SMB
attendu négatif). Aucune ne peut renverser le verdict : la \textbf{MDE vaut 7,8 \%/an}, deux ordres de grandeur
au-dessus. \emph{Le notebook 08 fait déjà Carhart-4 avec Newey-West : sur ce point le niveau membre est en retard
sur le niveau titre.}}

% ═══════════════════════════════════════════════════════════
\rem{A2}{La winsorisation --- elle ampute la queue qu'on cherche}
\cit{avoir avec et sans winsor : s'il y justement de l'inside trading, ca peut aussi être pour capturer ces micro evènements qui ont un gros impact}

\trad{le clip $r \mapsto \max(-c, \min(c, r))$ protège des \textbf{erreurs de cotation}. Mais il ne distingue pas une
erreur d'un \textbf{événement} : il coupe les deux. Il faut un critère qui les sépare.}

\textbf{Ce que le clip coûte, exactement.} Sur la NAV composée, l'amputation est multiplicative :
\[
\frac{\mathrm{NAV}^{\text{clip}}(T)}{\mathrm{NAV}(T)} \;=\; \prod_{t \,:\, r_t > c} \frac{1+c}{1+r_t}
\;\times\!\!\prod_{t \,:\, r_t < -c} \frac{1-c}{1+r_t}
\]
Le premier produit est $< 1$ et le second $> 1$ : \textbf{le clip est asymétrique par construction}, car
$r_t \in (-1, +\infty)$ --- borné en bas, pas en haut.

\textbf{Le critère qui sépare} : une erreur de cotation \textbf{se retourne} le lendemain, un événement \textbf{persiste}.
\[
\Lambda_{x,\varepsilon}(i,t) = \mathbf{1}\Bigl\{\, \underbrace{|r_i(t)| > x}_{x\,=\,50\,\%}
\;\wedge\;
\underbrace{\bigl|\,(1+r_i(t))(1+r_i(t{+}1)) - 1\,\bigr| < \varepsilon}_{\varepsilon\,=\,10\,\%\ :\ \text{l'aller-retour}} \,\Bigr\}
\]
On ne neutralise que les $(i,t)$ tels que $\Lambda = 1$. \textbf{Nichage} : $\varepsilon \to \infty$ rend $\Lambda$ vrai pour
tout saut $> x$ et redonne le clip brut.

\mes{sur \textbf{10\,340\,382} jours-titre, \textbf{3\,316} sauts dépassent $\pm 50\,\%$ --- dont \textbf{64 \% à la hausse}.
Se retournent : 31,3 \% des hausses, 14,8 \% des baisses. \textbf{2\,478 sauts persistent, soit 74,8 \% : le clip actuel
détruit trois vrais événements sur quatre}, et en majorité des hausses.}

{\small\textbf{Le test qui tranche} (celui qu'il demande implicitement) : soit $p_{\text{in}}$ la fréquence de clip dans
$[\delta_k, \delta_k{+}H]$ et $p_{\text{out}}$ hors de ces fenêtres. Si $p_{\text{in}} > p_{\text{out}}$ significativement,
le clip mange le signal lui-même. \textbf{Protocole} : un run primaire pré-enregistré (filtre $\Lambda$), les autres
réglages en sensibilité déclarée \textbf{non revendicable} --- sinon trois essais de plus au compteur.}

% ═══════════════════════════════════════════════════════════
\rem{A4}{Les seuils d'éligibilité --- pondérer par la précision au lieu de trancher}
\cit{assumption à relax}

\trad{$n_{\text{trades}} \ge 10$ et $n_{\text{jours}} \ge 126$ sont des \textbf{indicatrices} là où le problème est une
question de \textbf{précision}. Une porte binaire jette de l'information et ne protège pas : les deux
\og significatifs \fg{} les plus suspects la franchissent (Fetterman $t_{\mathrm{appr}}\,2{,}38$ sur 10 trades ;
Ruiz $t_{\mathrm{appr}}\,2{,}40$ sur 10 trades et \textbf{0 \% d'achats gagnants}).}

\textbf{Le modèle hiérarchique.} Chaque membre a un $\alpha$ propre, observé avec une erreur d'autant plus grande
que son historique est court :
\[
\alpha_m \sim \mathcal{N}(\mu, \tau^2),
\qquad
\hat\alpha_m \mid \alpha_m \sim \mathcal{N}(\alpha_m,\, s_m^2)
\]
\[
\Longrightarrow\qquad
\mathbb{E}[\alpha_m \mid \hat\alpha_m] = (1 - B_m)\,\hat\alpha_m + B_m\,\mu,
\qquad
B_m = \frac{s_m^2}{\tau^2 + s_m^2}
\]
\textbf{Estimation par les moments} (rien à calibrer à la main) :
\[
\hat\mu = \frac{\sum_m \hat\alpha_m / (\hat\tau^2 + s_m^2)}{\sum_m 1/(\hat\tau^2 + s_m^2)},
\qquad
\hat\tau^2 = \max\Bigl\{0,\;\; \widehat{\mathrm{Var}}(\hat\alpha) - \overline{s_m^2} \Bigr\}
\]

\prop{8 (le rétrécissement fait le travail de la porte, en continu)}{$B_m \to 1$ quand $s_m$ est grand (historique
court) : l'estimation est ramenée vers $\mu$. $B_m \to 0$ quand $s_m$ est petit : un membre à 700 trades ne bouge
presque pas. Comme $s_m^2 \propto 1/n_{\text{eff}}(m)$, \textbf{la précision se mesure en décisions} (\S A1), pas en
lignes : le rétrécissement et le décompte de rafale sont le même geste vu de deux côtés.
\textbf{Nichage} : $\tau^2 \to \infty$ donne $B_m \to 0$, aucun rétrécissement --- avec la porte, c'est le classement actuel.}

\mes{le seuil en décisions au lieu de lignes fait passer les membres jugeables de \textbf{264} ($n \ge 10$) à
\textbf{207} ($n_{\text{eff}} \ge 10$). \textbf{Fetterman et Ruiz doivent tomber} du classement a posteriori : c'est le test
de validité de la méthode, écrit avant de la lancer.}

{\small La sensibilité se présente en \textbf{carte} $(n_{\min}, d_{\min})$ avec le verdict en chaque point : on la lit
\og le verdict est-il stable partout ? \fg{}, \textbf{jamais} \og quel est le meilleur point \fg{}.
\emph{Ancrage : Blonien, Crane \& Crotty (2025) appliquent ce modèle hiérarchique à ces mêmes données.}}

% ═══════════════════════════════════════════════════════════
\vspace{2mm}
\begin{mdframed}[linecolor=rouge, backgroundcolor=rouge!5, linewidth=0.7pt, roundcorner=3pt]
\textbf{Le budget d'essais.} Ces adaptations ajoutent des spécifications, donc des occasions de trouver un faux
positif. Le compteur est à \textbf{160 essais}, avec $\mathbb{E}[\max t \mid H_0] = 2{,}69$ et un Deflated Sharpe de 0,11.
Discipline : \textbf{un seul run primaire pré-enregistré par remarque}, tout le reste en sensibilité déclarée
\textbf{non revendicable}, et $\mathbb{E}[\max t \mid H_0]$ recalculé pour le nouveau $N$ et publié avec les résultats.

\vspace{1mm}
\textbf{Ce qui ne changera pas : le verdict.} La MDE du dispositif est de \textbf{7,8 \%/an} sur onze années ; les
corrections ci-dessus déplacent $\alpha$ de quelques dixièmes de point. Elles ne rendent pas la stratégie
gagnante --- elles rendent la mesure \textbf{défendable}, et c'est ce qui était demandé.
\end{mdframed}

\end{document}
